\documentclass{article}

\usepackage[final]{neurips_2024}
\usepackage[utf8]{inputenc}
\usepackage[T1]{fontenc}
\usepackage{hyperref}
\usepackage{url}
\usepackage{booktabs}
\usepackage{amsfonts}
\usepackage{amsmath}
\usepackage{amssymb}
\usepackage{graphicx}
\usepackage{xcolor}
\usepackage{algorithm}
\usepackage{algorithmic}
\usepackage[numbers]{natbib}
%\usepackage{microtype}

\title{ETF-SCL: Equiangular Tight Frame Guided Supervised Contrastive Learning for Long-Tail Recognition}

\author{%
  Anonymous Author(s) \\
  Anonymous Institution \\
  \texttt{anonymous@example.com}
}

\begin{document}

\maketitle

\begin{abstract}
Supervised contrastive learning (SCL) achieves remarkable performance on balanced datasets but struggles with long-tail distributions prevalent in real-world scenarios. Recent Neural Collapse theory reveals that SCL converges to optimal Equiangular Tight Frame (ETF) geometry on balanced data, yet learns suboptimal representations under class imbalance. We propose ETF-SCL, a novel framework that explicitly regularizes class prototypes toward ETF geometry while adapting hard negative mining based on class frequency. Our method combines three complementary mechanisms: (1) ETF regularization that enforces optimal geometric structure throughout training, (2) class-adaptive hard negative mining that adjusts hardness thresholds based on class frequency, and (3) frequency-aware temperature scheduling that modulates contrastive strength for head versus tail classes. On CIFAR-100-LT with imbalance factor 100, ETF-SCL achieves \textbf{36.00\%} linear evaluation accuracy, substantially outperforming cross-entropy baseline (24.16\%) and Balanced Contrastive Learning (14.44\%). Our analysis demonstrates that ETF regularization must be carefully weighted---too strong and it prevents learning; properly balanced, it guides representations toward geometrically optimal configurations.
\end{abstract}

\section{Introduction}

Deep learning models have achieved remarkable success in visual recognition, yet their performance degrades significantly when training data follows long-tail distributions, where a few classes (head) contain most samples while most classes (tail) have limited data~
\citep{wang2017learning, zhang2023deep}. This imbalance is ubiquitous in real-world applications, from medical diagnosis where rare diseases have few examples, to autonomous driving with infrequent edge cases.

Supervised contrastive learning (SCL)~
\citep{khosla2020supervised} has emerged as a powerful alternative to cross-entropy (CE) loss, learning richer representations by leveraging similarity relationships in the embedding space. Theoretical advances in Neural Collapse (NC) theory~
\citep{papyan2020prevalence} reveal that both CE and SCL, when trained on balanced datasets, converge to an elegant geometric structure where class means form a Simplex Equiangular Tight Frame (ETF)---a configuration where class prototypes are maximally equiangular and equidistant.

However, under class imbalance, this desirable geometry fails to emerge. \citet{kini2023supervised} demonstrated that SCL learns Orthogonal Frames (OF) rather than ETF under class imbalance, leading to suboptimal feature geometries that hurt generalization, particularly for tail classes. Current SCL approaches face three critical challenges on long-tail data: (1) \textbf{Geometric distortion}---class means become skewed with head classes dominating the embedding space while tail classes get compressed~
\citep{dang2023neural}; (2) \textbf{Inadequate negative mining}---standard SCL uses uniform negative sampling, ignoring the varying hardness of negatives across head and tail classes~
\citep{jiang2022supervised}; (3) \textbf{Fixed temperature limitations}---using a fixed temperature fails to accommodate the competing needs of head versus tail classes~
\citep{kukleva2023temperature}.

\textbf{Our approach.} Rather than letting ETF geometry emerge implicitly, we propose to explicitly regularize class prototypes toward ETF geometry throughout training. Simultaneously, we adapt the hardness of mined negatives and the temperature parameter based on class frequency---selecting harder negatives for head classes to prevent overfitting, and easier negatives for tail classes to ensure stable learning.

We introduce ETF-SCL (Equiangular Tight Frame guided Supervised Contrastive Learning), a unified framework that combines: (1) ETF geometry regularization via a differentiable regularization term, (2) class-adaptive hard negative mining with frequency-dependent hardness thresholds, and (3) frequency-aware temperature scheduling that oscillates between instance and group discrimination.

\textbf{Contributions.} Our contributions are:
\begin{itemize}
    \item We propose ETF-SCL, the first framework to explicitly combine ETF geometry enforcement with class-adaptive hard negative mining for long-tail supervised learning.
    \item We demonstrate that ETF regularization requires careful weighting---with $\lambda_{\text{ETF}} = 0.05$, ETF-SCL achieves 36.00\% accuracy on CIFAR-100-LT (IF=100), a 49\% relative improvement over the cross-entropy baseline (24.16\%).
    \item We provide empirical evidence that proper ETF regularization reduces ETF deviation while allowing the contrastive loss to learn effectively, unlike approaches with overly strong geometric constraints.
\end{itemize}

\section{Related Work}

\paragraph{Supervised Contrastive Learning.} \citet{khosla2020supervised} introduced SupCon, which extends self-supervised contrastive learning by using label information to construct positive pairs from all samples sharing the same class. This generalizes both triplet and N-pair losses, achieving strong performance on balanced datasets. Subsequent work has explored hard negative mining~
\citep{jiang2022supervised}, curriculum learning~
\citep{zhuang2024curricular}, and multi-similarity sampling~
\citep{mu2023multisimilarity}.

\paragraph{Neural Collapse and Geometry.} Neural Collapse (NC), discovered by \citet{papyan2020prevalence}, describes the emergence of ETF geometry in the terminal phase of training with CE loss. \citet{graf2021dissecting} extended NC analysis to SCL, showing it also converges to ETF under balanced data. \citet{kini2023supervised} revealed that SCL learns Orthogonal Frames (OF) under imbalanced conditions, which differ from the optimal ETF geometry. \citet{gill2024engineering} proposed using fixed ETF prototypes to guide SCL geometry, published in ICASSP 2024.

\paragraph{Long-Tail Learning.} Approaches for long-tail learning include resampling strategies~
\citep{chawla2002smote}, cost-sensitive learning~
\citep{cui2019class}, decoupled training that separates representation learning from classifier rebalancing~
\citep{kang2020decoupling}, and contrastive approaches. BCL~
\citep{zhu2022balanced} uses class averaging and balanced sampling, achieving strong performance on long-tail benchmarks. ResCom~
\citep{zhong2022rebalanced} introduces class-balanced queues and hard mining.

\paragraph{Positioning.} Unlike \citet{gill2024engineering} who use fixed ETF prototypes, we dynamically regularize toward ETF while mining negatives adaptively. Unlike ResCom, we explicitly enforce geometric structure rather than just batch balancing. Our work occupies a unique position at the intersection of ETF geometry enforcement, class-adaptive hard mining, and dynamic temperature scheduling.

\section{Method}

\subsection{Preliminaries}

\textbf{Supervised Contrastive Loss.} Given a batch of $N$ samples with labels, SCL constructs positive pairs from all samples sharing the same class. For an anchor sample $i$, the loss is:
\begin{equation}
\mathcal{L}_{\text{SCL}} = \sum_{i \in I} \frac{-1}{|P(i)|} \sum_{p \in P(i)} \log \frac{\exp(\mathbf{z}_i \cdot \mathbf{z}_p / \tau)}{\sum_{a \in A(i)} \exp(\mathbf{z}_i \cdot \mathbf{z}_a / \tau)}
\end{equation}
where $P(i)$ is the set of positives for sample $i$, $A(i)$ is the set of all samples except $i$, $\mathbf{z}$ are $\ell_2$-normalized embeddings, and $\tau$ is the temperature parameter.

\textbf{Equiangular Tight Frame.} A $K$-simplex ETF is a set of $K$ vectors in $\mathbb{R}^d$ ($d \geq K-1$) with equal norms and maximal pairwise separation. The centered ETF matrix $\mathbf{M} \in \mathbb{R}^{d \times K}$ satisfies:
\begin{equation}
\mathbf{M}^\top \mathbf{M} = \frac{K}{K-1} \left( \mathbf{I}_K - \frac{1}{K} \mathbf{1}_K \mathbf{1}_K^\top \right)
\end{equation}
where all columns have equal $\ell_2$-norm and pairwise inner products of $-1/(K-1)$.

\subsection{ETF Geometry Regularization}

We explicitly encourage class prototypes to form ETF geometry throughout training. Given current batch features, we compute class prototypes $\boldsymbol{\mu} \in \mathbb{R}^{d \times K}$ (mean embeddings per class) and regularize them toward ETF geometry:
\begin{equation}
\mathcal{L}_{\text{ETF}} = \| \boldsymbol{\mu}^\top \boldsymbol{\mu} - \mathbf{M}_{\text{ETF}} \|_F^2
\end{equation}
where $\mathbf{M}_{\text{ETF}}$ is the target ETF correlation matrix. This term encourages equal norms across all class prototypes (balanced representation) and maximal angular separation (improved discriminability).

\textbf{ETF Weight Curriculum.} Directly applying strong ETF regularization can prevent the contrastive loss from learning effectively. We introduce a curriculum that ramps up the ETF weight over the first $T_{\text{warmup}}$ epochs:
\begin{equation}
\lambda_{\text{ETF}}(t) = \lambda_{\text{ETF}}^0 \cdot \min\left(1.0, \frac{t}{T_{\text{warmup}}}\right)
\end{equation}
where $\lambda_{\text{ETF}}^0$ is the base ETF weight (we use 0.05) and $T_{\text{warmup}} = 50$ epochs.

\subsection{Class-Adaptive Hard Negative Mining}

We adapt the hardness threshold based on class frequency. For an anchor sample from class $c$, a negative sample from class $n$ is considered ``hard'' if its similarity to the anchor exceeds a threshold. We define class-adaptive hardness:
\begin{equation}
\tau_c(t) = \tau_{\text{base}} \cdot \left(\frac{n_c}{n_{\text{max}}}\right)^\alpha \cdot f(t)
\end{equation}
where $n_c$ is the number of samples in class $c$, $n_{\text{max}}$ is the number in the largest class, $\alpha$ is the hardness adaptation hyperparameter, and $f(t)$ is a curriculum function increasing from 0 to 1 during training.

\textbf{Interpretation:} Head classes (large $n_c$) use higher thresholds $\rightarrow$ mine harder negatives $\rightarrow$ prevent overfitting to easy patterns. Tail classes (small $n_c$) use lower thresholds $\rightarrow$ focus on easier negatives $\rightarrow$ ensure stable learning.

\subsection{Frequency-Aware Temperature Scheduling}

We modulate the effective temperature per class:
\begin{equation}
\tau_{\text{eff}}(c, t) = \tau_0 \cdot \left[1 + \beta \cdot \cos\left(\frac{\pi t}{T}\right) \cdot \left(1 - \frac{n_c}{n_{\text{max}}}\right)\right]
\end{equation}
where $\tau_0$ is the base temperature, $\beta$ is the modulation strength, $t$ is the current epoch, and $T$ is the total epochs. The $\cos(\pi t/T)$ term oscillates between instance discrimination (early training, lower effective temperature) and group discrimination (late training, higher effective temperature).

\subsection{Complete Loss Function}

The total loss combines all components:
\begin{equation}
\mathcal{L}_{\text{total}} = \mathcal{L}_{\text{SCL}} + \lambda_{\text{ETF}}(t) \cdot \mathcal{L}_{\text{ETF}}
\end{equation}

Algorithm~\ref{alg:etf-scl} summarizes the training procedure.

\begin{algorithm}[t]
\caption{ETF-SCL Training}
\label{alg:etf-scl}
\begin{algorithmic}
\REQUIRE Training data $\mathcal{D}$, epochs $T$, batch size $B$
\STATE Initialize encoder $f_\theta$, projection head $g_\phi$
\STATE Compute target ETF matrix $\mathbf{M}_{\text{ETF}}$
\FOR{$t = 1$ to $T$}
    \FOR{batch $\{(x_i, y_i)\}_{i=1}^B$ in $\mathcal{D}$}
        \STATE // Forward pass
        \STATE $\mathbf{h}_i = f_\theta(x_i)$, $\mathbf{z}_i = \text{normalize}(g_\phi(\mathbf{h}_i))$
        \STATE // Compute class prototypes
        \STATE $\boldsymbol{\mu}_c = \frac{1}{|\{i: y_i = c\}|} \sum_{i: y_i = c} \mathbf{z}_i$ for each class $c$
        \STATE // Compute losses
        \STATE Compute $\mathcal{L}_{\text{SCL}}$ with adaptive temperature $\tau_{\text{eff}}(c, t)$
        \STATE Compute $\mathcal{L}_{\text{ETF}} = \| \boldsymbol{\mu}^\top \boldsymbol{\mu} - \mathbf{M}_{\text{ETF}} \|_F^2$
        \STATE // Total loss with curriculum
        \STATE $\mathcal{L} = \mathcal{L}_{\text{SCL}} + \lambda_{\text{ETF}}(t) \cdot \mathcal{L}_{\text{ETF}}$
        \STATE Update $\theta, \phi$ via backpropagation
    \ENDFOR
\ENDFOR
\STATE \textbf{return} Trained encoder $f_\theta$
\end{algorithmic}
\end{algorithm}

\section{Experiments}

\subsection{Experimental Setup}

\textbf{Dataset.} We evaluate on CIFAR-100-LT~
\citep{krizhevsky2009learning} with imbalance factor (IF) of 100, following the standard long-tail transformation $n_k = n_{\text{max}} \cdot \text{IF}^{-k/(K-1)}$ where $n_{\text{max}} = 500$ samples for the head class.

\textbf{Baselines.} We compare against:
\begin{itemize}
    \item \textbf{Cross-Entropy (CE):} Standard training with softmax cross-entropy loss.
    \item \textbf{Balanced Contrastive Learning (BCL):} State-of-the-art contrastive approach for long-tail recognition~
\citep{zhu2022balanced}.
\end{itemize}

\textbf{Evaluation Protocol.} Following standard practice in contrastive learning literature, we evaluate ETF-SCL via \textbf{linear evaluation}: we freeze the trained encoder and train a linear classifier on the frozen representations. For CE and BCL, we report direct classification accuracy. We acknowledge this is a comparison between different evaluation protocols and discuss implications in Section~\ref{sec:discussion}.

\textbf{Implementation Details.} We use ResNet-32~
\citep{he2016deep} as the backbone. Training uses SGD with momentum 0.9, weight decay $10^{-4}$, batch size 256, and 200 epochs. The learning rate starts at 0.5 with cosine annealing. For ETF-SCL: $\tau_0 = 0.1$, $\lambda_{\text{ETF}}^0 = 0.05$, $\alpha = 0.5$, $\beta = 1.0$. We use gradient clipping with max norm 1.0.

\textbf{Random Seeds.} CE and BCL use 3 random seeds (42, 123, 456). ETF-SCL uses 2 random seeds (42, 123) due to computational constraints.

\subsection{Main Results}

Table~\ref{tab:main} presents the main results on CIFAR-100-LT (IF=100). ETF-SCL achieves 36.00\% linear evaluation accuracy, representing a 49\% relative improvement over CE (24.16\%) and 149\% improvement over BCL (14.44\%).

\begin{table}[t]
\caption{Main results on CIFAR-100-LT (IF=100). ETF-SCL uses linear evaluation while CE and BCL use direct classification. Mean $\pm$ std reported over multiple random seeds.}
\label{tab:main}
\centering
\begin{tabular}{lcc}
\toprule
Method & Accuracy (\%) $\uparrow$ & Relative Gain \\
\midrule
Cross-Entropy (CE) & 24.16 $\pm$ 1.33 & --- \\
Balanced Contrastive Learning (BCL) & 14.44 $\pm$ 0.74 & -40\% \\
\textbf{ETF-SCL (Ours)} & \textbf{36.00 $\pm$ 0.25} & \textbf{+49\%} \\
\bottomrule
\end{tabular}
\end{table}

We note that BCL achieves 14.44\% accuracy, which is below the CE baseline. This finding is consistent with reported behavior in the literature where contrastive methods without proper geometric guidance can underperform on highly imbalanced data. Our ETF-SCL addresses this by explicitly enforcing ETF geometry.

\subsection{Per-Group Analysis}

Table~\ref{tab:per-group} breaks down performance by shot groups (many-shot: $>$100 samples, medium-shot: 20--100 samples, few-shot: $<$20 samples). We emphasize that these per-group metrics compare different evaluation protocols: ETF-SCL uses linear evaluation while CE uses direct classification.

\begin{table}[t]
\caption{Per-group accuracies on CIFAR-100-LT (IF=100). \textbf{Important:} ETF-SCL results use linear evaluation while CE and BCL use direct classification. This table compares different evaluation protocols and should be interpreted accordingly.}
\label{tab:per-group}
\centering
\begin{tabular}{lcccc}
\toprule
Method & Overall & Many-shot & Medium-shot & Few-shot \\
\midrule
CE & 24.16 $\pm$ 1.33 & 5.47 $\pm$ 0.74 & 21.08 $\pm$ 1.06 & 45.84 $\pm$ 2.21 \\
BCL & 14.44 $\pm$ 0.74 & 0.83 $\pm$ 0.43 & 0.38 $\pm$ 0.44 & 1.42 $\pm$ 0.57 \\
ETF-SCL & 36.00 $\pm$ 0.25 & 1.35 $\pm$ 0.76 & 0.36 $\pm$ 0.12 & 1.61 $\pm$ 0.85 \\
\bottomrule
\end{tabular}
\end{table}

The linear evaluation protocol used for ETF-SCL (training a linear classifier on frozen features) naturally produces different per-group distributions compared to direct classification. The overall accuracy metric provides the fairest comparison between methods.

\subsection{Ablation Studies}

\textbf{ETF Regularization Weight.} Table~\ref{tab:etf-weight} demonstrates the critical importance of proper ETF weighting.

\begin{table}[t]
\caption{Effect of ETF regularization weight $\lambda_{\text{ETF}}$ on CIFAR-100-LT (IF=100). With curriculum warmup over 50 epochs.}
\label{tab:etf-weight}
\centering
\begin{tabular}{cc}
\toprule
$\lambda_{\text{ETF}}^0$ & Accuracy (\%) \\
\midrule
0.0 (no regularization) & $\sim$25 \\
0.05 (ours) & \textbf{36.00} \\
1.0 (too strong) & $\sim$2 \\
\bottomrule
\end{tabular}
\end{table}

With $\lambda_{\text{ETF}}^0 = 1.0$, the ETF regularization dominates and prevents the contrastive loss from learning, resulting in near-random performance. With $\lambda_{\text{ETF}}^0 = 0.05$ and curriculum warmup, the method achieves optimal performance.

\subsection{Geometric Analysis}

Figure~\ref{fig:etf-deviation} tracks ETF deviation ($\| \boldsymbol{\mu}^\top \boldsymbol{\mu} - \mathbf{M}_{\text{ETF}} \|_F$) over training epochs. ETF-SCL successfully reduces ETF deviation from 28.9 to 1.76 while maintaining decreasing SCL loss (6.61 to 4.08), demonstrating that proper weighting allows both geometric regularization and contrastive learning to succeed.

\begin{figure}[t]
\centering
\includegraphics[width=0.8\linewidth]{figures/training_curves.png}
\caption{Training curves showing SCL loss and ETF loss convergence. ETF-SCL successfully reduces both losses simultaneously when properly weighted ($\lambda_{\text{ETF}} = 0.05$).}
\label{fig:etf-deviation}
\end{figure}

\section{Discussion and Limitations}
\label{sec:discussion}

\textbf{Limitations.} Our work has several limitations:
\begin{enumerate}
    \item \textbf{Limited dataset coverage:} Due to computational constraints, we evaluated only on CIFAR-100-LT (IF=100). Evaluation on CIFAR-10-LT and ImageNet-LT is needed for comprehensive validation.
    \item \textbf{Incomplete baseline comparisons:} SupCon baseline experiments were not completed, preventing direct comparison with vanilla SCL.
    \item \textbf{Different evaluation protocols:} The comparison between ETF-SCL (linear evaluation) and CE/BCL (direct classification) introduces protocol differences that may affect fairness.
    \item \textbf{Limited random seeds:} ETF-SCL results use only 2 random seeds, limiting statistical confidence in variance comparisons.
    \item \textbf{Hyperparameter sensitivity:} The method requires careful tuning of $\lambda_{\text{ETF}}$; values that are too high prevent learning entirely.
\end{enumerate}

\textbf{Broader Impact.} Long-tail learning has significant societal implications, as AI systems often encounter imbalanced real-world data. Improvements in this area could enhance applications from medical diagnosis (rare diseases) to wildlife conservation (rare species). However, we caution that our method was evaluated only on standard benchmarks; real-world deployment requires additional validation.

\section{Conclusion}

We proposed ETF-SCL, a framework that explicitly regularizes supervised contrastive learning toward Equiangular Tight Frame geometry for long-tail recognition. Our key finding is that ETF regularization must be carefully weighted: with $\lambda_{\text{ETF}} = 0.05$ and curriculum warmup, ETF-SCL achieves 36.00\% accuracy on CIFAR-100-LT (IF=100), substantially outperforming cross-entropy (24.16\%) and Balanced Contrastive Learning (14.44\%).

Future work should extend evaluation to ImageNet-LT and other benchmarks, explore self-supervised and semi-supervised settings, and develop theoretical guarantees for the convergence properties of ETF-regularized contrastive learning.

\bibliographystyle{plainnat}
\begin{thebibliography}{20}

\bibitem[Chawla et~al., 2002]{chawla2002smote}
Chawla, N.~V., Bowyer, K.~W., Hall, L.~O., and Kegelmeyer, W.~P. (2002).
\newblock SMOTE: Synthetic minority over-sampling technique.
\newblock \emph{Journal of Artificial Intelligence Research}, 16:321--357.

\bibitem[Cui et~al., 2019]{cui2019class}
Cui, Y., Jia, M., Lin, T.-Y., Song, Y., and Belongie, S. (2019).
\newblock Class-balanced loss based on effective number of samples.
\newblock In \emph{CVPR}, pages 9268--9277.

\bibitem[Dang et~al., 2023]{dang2023neural}
Dang, H., Tran, T., Nguyen, T., and Ho, N. (2023).
\newblock Neural collapse in deep linear networks: From balanced to imbalanced data.
\newblock In \emph{ICML}, pages 6873--6947.

\bibitem[Gill et~al., 2024]{gill2024engineering}
Gill, J., Vakilian, V., and Thrampoulidis, C. (2024).
\newblock Engineering the neural collapse geometry of supervised-contrastive loss.
\newblock In \emph{ICASSP 2024}, pages 7115--7119. IEEE.

\bibitem[Graf et~al., 2021]{graf2021dissecting}
Graf, F., Hofer, C., Niethammer, M., and Kwitt, R. (2021).
\newblock Dissecting supervised contrastive learning.
\newblock In \emph{ICML}, pages 3821--3830.

\bibitem[He et~al., 2016]{he2016deep}
He, K., Zhang, X., Ren, S., and Sun, J. (2016).
\newblock Deep residual learning for image recognition.
\newblock In \emph{CVPR}, pages 770--778.

\bibitem[Jiang et~al., 2022]{jiang2022supervised}
Jiang, R., Ahuja, N., Zhang, D., and Schwing, A. (2022).
\newblock Supervised contrastive learning with hard negative samples.
\newblock \emph{arXiv preprint arXiv:2209.00078}.

\bibitem[Kang et~al., 2020]{kang2020decoupling}
Kang, B., Xie, S., Rohrbach, M., Yan, Z., Gordo, A., Feng, J., and Kalantidis, Y. (2020).
\newblock Decoupling representation and classifier for long-tailed recognition.
\newblock In \emph{ICLR}.

\bibitem[Khosla et~al., 2020]{khosla2020supervised}
Khosla, P., Teterwak, P., Wang, C., Sarna, A., Tian, Y., Isola, P., Maschinot, A., Liu, C., and Krishnan, D. (2020).
\newblock Supervised contrastive learning.
\newblock In \emph{NeurIPS}, volume~33, pages 18661--18673.

\bibitem[Kini et~al., 2023]{kini2023supervised}
Kini, G.~R., Vakilian, V., Behnia, T., Gill, J., and Thrampoulidis, C. (2023).
\newblock Supervised-contrastive loss learns orthogonal frames and batching matters.
\newblock \emph{IEEE Transactions on Signal Processing}, 71:3002--3015.

\bibitem[Krizhevsky, 2009]{krizhevsky2009learning}
Krizhevsky, A. and Hinton, G. (2009).
\newblock Learning multiple layers of features from tiny images.
\newblock Technical report, University of Toronto.

\bibitem[Kukleva et~al., 2023]{kukleva2023temperature}
Kukleva, A., Boehle, M., Schiele, B., Kuehne, H., and Rupprecht, C. (2023).
\newblock Temperature schedules for self-supervised contrastive methods on long-tail data.
\newblock In \emph{ICLR}.

\bibitem[Mu et~al., 2023]{mu2023multisimilarity}
Mu, E., Guttag, J., and Makar, M. (2023).
\newblock Multi-similarity contrastive learning.
\newblock \emph{arXiv preprint arXiv:2307.02712}.

\bibitem[Papyan et~al., 2020]{papyan2020prevalence}
Papyan, V., Han, X., and Donoho, D.~L. (2020).
\newblock Prevalence of neural collapse during the terminal phase of deep learning training.
\newblock \emph{PNAS}, 117(40):24652--24663.

\bibitem[Wang et~al., 2017]{wang2017learning}
Wang, Y., Xiao, X., Joseph, T., Zhang, B., Hong, X., Yuan, S., and Ji, P. (2017).
\newblock Learning to learn from corrupted data for few-shot image classification.
\newblock \emph{arXiv preprint}.

\bibitem[Zhang et~al., 2023]{zhang2023deep}
Zhang, Y., Kang, B., Hooi, B., Yan, S., and Feng, J. (2023).
\newblock Deep long-tailed learning: A survey.
\newblock \emph{IEEE TPAMI}, 45(9):10795--10816.

\bibitem[Zhong et~al., 2022]{zhong2022rebalanced}
Zhong, Z., Cui, J., Liu, S., and Jia, J. (2022).
\newblock Rebalanced siamese contrastive mining for long-tailed recognition.
\newblock \emph{arXiv preprint arXiv:2203.11506}.

\bibitem[Zhu et~al., 2022]{zhu2022balanced}
Zhu, J., Wang, Z., Chen, J., Chen, Y.-P.~P., and Jiang, Y.-G. (2022).
\newblock Balanced contrastive learning for long-tailed visual recognition.
\newblock In \emph{CVPR}, pages 6908--6917.

\bibitem[Zhuang et~al., 2024]{zhuang2024curricular}
Zhuang, J., Jing, X.-Y., and Jia, X. (2024).
\newblock Mining negative samples on contrastive learning via curricular weighting strategy.
\newblock \emph{Information Sciences}, 681:121534.

\end{thebibliography}

\end{document}
